\documentclass[12pt,a4paper]{article}
\usepackage[UTF8]{ctex}
\usepackage{amsmath,amssymb,amsthm}
\usepackage{geometry}

\geometry{margin=2.5cm}

\title{混合运动-力控制器的速度控制版本}
\author{第11章补充说明}
\date{}

\begin{document}

\maketitle

\section{问题提出}

原始混合运动-力控制器公式(11.61)：
\begin{equation}
\tau = J_b^T(\theta)\Bigg[P(\theta)\left(\tilde{\Lambda}(\theta)\frac{d}{dt}([Ad_{X^{-1}X_d}]V_d) + K_p X_e + K_i\int_0^t X_e(t)dt + K_d V_e\right) + (I - P(\theta))\left(F_d + K_{fp}F_e + K_{fi}\int_0^t F_e(t)dt\right) + \tilde{\eta}(\theta, V_b)\Bigg] \tag{11.61}
\end{equation}

\textbf{问题}：如果控制的是末端速度而不是加速度，公式应该怎么写？

\section{加速度控制 vs 速度控制}

\subsection{加速度控制（原始公式）}

\textbf{特点}：
\begin{itemize}
\item 控制器输出期望加速度$\dot{V}_d$
\item 需要计算$\frac{d}{dt}([Ad_{X^{-1}X_d}]V_d)$
\item 使用质量矩阵$\tilde{\Lambda}(\theta)$将加速度转换为力
\item 适用于可以控制加速度的系统
\end{itemize}

\textbf{运动控制部分}：
\begin{equation}
F_{\text{motion}} = \tilde{\Lambda}(\theta)\frac{d}{dt}([Ad_{X^{-1}X_d}]V_d) + K_p X_e + K_i\int_0^t X_e(t)dt + K_d V_e
\end{equation}

\subsection{速度控制}

\textbf{特点}：
\begin{itemize}
\item 控制器直接输出期望速度$V_d$
\item 不需要计算加速度
\item 不需要质量矩阵（或使用简化的质量矩阵）
\item 适用于只能控制速度的系统
\end{itemize}

\textbf{运动控制部分}：
\begin{equation}
F_{\text{motion}} = K_p X_e + K_i\int_0^t X_e(t)dt + K_d V_e
\end{equation}

\section{速度控制版本的推导}

\subsection{基本思路}

\textbf{目标}：设计速度控制器，使得末端执行器速度$V_b$跟踪期望速度$V_d$。

\textbf{方法}：
\begin{enumerate}
\item 定义速度误差：$V_e = V_d - V_b$
\item 使用PID控制器产生控制力
\item 投影到运动控制子空间
\item 与力控制器组合
\end{enumerate}

\subsection{速度控制器的设计}

\textbf{速度误差}：
\begin{equation}
V_e = V_d - V_b
\end{equation}

其中：
\begin{itemize}
\item $V_d$：期望的末端执行器速度（在体坐标系$\{b\}$中）
\item $V_b$：实际的末端执行器速度（在体坐标系$\{b\}$中）
\end{itemize}

\textbf{位置误差}：
\begin{equation}
X_e = \log(X^{-1}X_d)
\end{equation}

其中：
\begin{itemize}
\item $X$：实际的末端执行器位姿
\item $X_d$：期望的末端执行器位姿
\item $X_e$：位置误差（在SE(3)中）
\end{itemize}

\textbf{速度控制器}：
\begin{equation}
F_{\text{motion}} = K_p X_e + K_i\int_0^t X_e(\tau)d\tau + K_d V_e
\end{equation}

其中：
\begin{itemize}
\item $K_p$：位置比例增益矩阵
\item $K_i$：位置积分增益矩阵
\item $K_d$：速度微分增益矩阵
\end{itemize}

\subsection{完整的速度控制公式}

\textbf{混合运动-力控制器（速度控制版本）}：
\begin{equation}
\tau = J_b^T(\theta)\Bigg[P(\theta)\left(K_p X_e + K_i\int_0^t X_e(\tau)d\tau + K_d V_e\right) + (I - P(\theta))\left(F_d + K_{fp}F_e + K_{fi}\int_0^t F_e(\tau)d\tau\right) + \tilde{\eta}(\theta, V_b)\Bigg] \tag{11.61-velocity}
\end{equation}

\textbf{与原始公式的对比}：

\textbf{原始公式（加速度控制）}：
\begin{equation}
F_{\text{motion}} = \tilde{\Lambda}(\theta)\frac{d}{dt}([Ad_{X^{-1}X_d}]V_d) + K_p X_e + K_i\int_0^t X_e(t)dt + K_d V_e
\end{equation}

\textbf{速度控制版本}：
\begin{equation}
F_{\text{motion}} = K_p X_e + K_i\int_0^t X_e(\tau)d\tau + K_d V_e
\end{equation}

\textbf{主要区别}：
\begin{itemize}
\item \textbf{移除了加速度项}：$\tilde{\Lambda}(\theta)\frac{d}{dt}([Ad_{X^{-1}X_d}]V_d)$
\item \textbf{保留了PID项}：$K_p X_e + K_i\int X_e dt + K_d V_e$
\item \textbf{力控制部分不变}：$(I-P)(F_d + K_{fp}F_e + K_{fi}\int F_e dt)$
\item \textbf{补偿项不变}：$\tilde{\eta}(\theta, V_b)$
\end{itemize}

\section{详细解释}

\subsection{各项的物理意义}

\textbf{运动控制部分}：$P(\theta)(K_p X_e + K_i\int X_e dt + K_d V_e)$

\begin{itemize}
\item \textbf{$K_p X_e$}：位置误差的比例项
\begin{itemize}
\item 产生与位置误差成比例的力
\item 使系统向期望位置移动
\end{itemize}

\item \textbf{$K_i\int X_e dt$}：位置误差的积分项
\begin{itemize}
\item 消除稳态误差
\item 累积位置误差，产生持续的控制力
\end{itemize}

\item \textbf{$K_d V_e$}：速度误差的微分项
\begin{itemize}
\item 产生与速度误差成比例的阻尼力
\item 稳定系统，减少振荡
\end{itemize}

\item \textbf{$P(\theta)$}：投影到运动控制子空间
\begin{itemize}
\item 确保控制力只影响自由运动方向
\item 不影响约束方向
\end{itemize}
\end{itemize}

\textbf{力控制部分}：$(I-P(\theta))(F_d + K_{fp}F_e + K_{fi}\int F_e dt)$

\begin{itemize}
\item \textbf{$F_d$}：期望的力
\item \textbf{$K_{fp}F_e$}：力误差的比例项
\item \textbf{$K_{fi}\int F_e dt$}：力误差的积分项
\item \textbf{$(I-P(\theta))$}：投影到力控制子空间
\end{itemize}

\textbf{补偿项}：$\tilde{\eta}(\theta, V_b)$

\begin{itemize}
\item 科里奥利力、向心力和重力补偿
\item 消除这些干扰项的影响
\end{itemize}

\subsection{速度误差的定义}

\textbf{关键问题}：如何定义$V_e$？

\textbf{方法1：直接差值}
\begin{equation}
V_e = V_d - V_b
\end{equation}

其中$V_d$和$V_b$都在体坐标系$\{b\}$中表示。

\textbf{方法2：考虑位姿变换}
\begin{equation}
V_e = [Ad_{X^{-1}X_d}]V_d - V_b
\end{equation}

这考虑了期望位姿$X_d$和实际位姿$X$之间的变换。

\textbf{推荐}：使用方法2，因为它考虑了位姿误差的影响。

\section{速度控制 vs 加速度控制}

\subsection{适用场景}

\textbf{加速度控制}：
\begin{itemize}
\item 适用于可以控制加速度的系统
\item 需要精确的动力学模型
\item 响应更快，性能更好
\item 但实现更复杂
\end{itemize}

\textbf{速度控制}：
\begin{itemize}
\item 适用于只能控制速度的系统
\item 不需要精确的加速度模型
\item 实现更简单
\item 但响应可能较慢
\end{itemize}

\subsection{性能对比}

\textbf{加速度控制}：
\begin{itemize}
\item \textbf{优点}：
\begin{itemize}
\item 可以精确控制加速度
\item 响应快
\item 性能好
\end{itemize}
\item \textbf{缺点}：
\begin{itemize}
\item 需要计算$\frac{d}{dt}([Ad_{X^{-1}X_d}]V_d)$
\item 需要质量矩阵$\tilde{\Lambda}(\theta)$
\item 对模型误差敏感
\end{itemize}
\end{itemize}

\textbf{速度控制}：
\begin{itemize}
\item \textbf{优点}：
\begin{itemize}
\item 实现简单
\item 不需要加速度计算
\item 对模型误差不敏感
\end{itemize}
\item \textbf{缺点}：
\begin{itemize}
\item 响应可能较慢
\item 性能可能不如加速度控制
\item 需要较大的增益来补偿
\end{itemize}
\end{itemize}

\section{实际实现}

\subsection{离散时间实现}

\textbf{速度控制器的离散时间形式}：

\textbf{位置误差}：
\begin{equation}
X_e(k) = \log(X^{-1}(k)X_d(k))
\end{equation}

\textbf{速度误差}：
\begin{equation}
V_e(k) = [Ad_{X^{-1}(k)X_d(k)}]V_d(k) - V_b(k)
\end{equation}

\textbf{积分项}：
\begin{equation}
I(k+1) = I(k) + X_e(k)\Delta t
\end{equation}

\textbf{控制力}：
\begin{equation}
F_{\text{motion}}(k) = K_p X_e(k) + K_i I(k+1) + K_d V_e(k)
\end{equation}

\textbf{完整控制器}：
\begin{equation}
\tau(k) = J_b^T(\theta(k))\Bigg[P(\theta(k))F_{\text{motion}}(k) + (I - P(\theta(k)))F_{\text{force}}(k) + \tilde{\eta}(\theta(k), V_b(k))\Bigg]
\end{equation}

\subsection{参数选择}

\textbf{增益矩阵的选择}：

\begin{itemize}
\item \textbf{$K_p$}：位置比例增益
\begin{itemize}
\item 较大的$K_p$：响应快，但可能不稳定
\item 较小的$K_p$：响应慢，但更稳定
\end{itemize}

\item \textbf{$K_i$}：位置积分增益
\begin{itemize}
\item 消除稳态误差
\item 但可能导致超调和振荡
\item 需要添加积分器抗饱和
\end{itemize}

\item \textbf{$K_d$}：速度微分增益
\begin{itemize}
\item 提供阻尼，稳定系统
\item 减少超调和振荡
\end{itemize}
\end{itemize}

\section{总结}

\subsection{速度控制版本的公式}

\textbf{混合运动-力控制器（速度控制版本）}：
\begin{equation}
\tau = J_b^T(\theta)\Bigg[P(\theta)\left(K_p X_e + K_i\int_0^t X_e(\tau)d\tau + K_d V_e\right) + (I - P(\theta))\left(F_d + K_{fp}F_e + K_{fi}\int_0^t F_e(\tau)d\tau\right) + \tilde{\eta}(\theta, V_b)\Bigg]
\end{equation}

其中：
\begin{itemize}
\item $X_e = \log(X^{-1}X_d)$：位置误差
\item $V_e = [Ad_{X^{-1}X_d}]V_d - V_b$：速度误差（推荐）
\item 或$V_e = V_d - V_b$：速度误差（简化）
\item $P(\theta)$：投影到运动控制子空间
\item $I - P(\theta)$：投影到力控制子空间
\end{itemize}

\subsection{与原始公式的对比}

\begin{center}
\begin{tabular}{|l|l|}
\hline
\textbf{原始公式（加速度控制）} & \textbf{速度控制版本} \\
\hline
包含加速度项$\tilde{\Lambda}\frac{d}{dt}([Ad]V_d)$ & 不包含加速度项 \\
需要质量矩阵$\tilde{\Lambda}$ & 不需要质量矩阵 \\
需要计算加速度 & 只需要速度 \\
响应快，性能好 & 响应较慢，但实现简单 \\
\hline
\end{tabular}
\end{center}

\subsection{选择建议}

\begin{itemize}
\item \textbf{如果可以控制加速度}：使用原始公式(11.61)
\item \textbf{如果只能控制速度}：使用速度控制版本
\item \textbf{如果系统响应要求高}：优先考虑加速度控制
\item \textbf{如果实现简单性优先}：使用速度控制
\end{itemize}

\end{document}

